\section{다관절 로봇을 읽기 위한 최소 문법}
\label{sec:robotics-foundations}

앞 장에서는 질량--스프링--댐퍼 하나만 보아도 힘, 위치, 속도, 에너지가 어떻게 얽히는지 꽤 많은 것을 설명할 수 있었다.
하지만 로봇 팔로 넘어오면 상황이 한 단계 복잡해진다.
사람은 보통 로봇 손끝이 어디에 있는지, 벽을 얼마나 눌렀는지, 물체를 어느 방향으로 밀었는지를 보고 싶어 한다.
반면 실제 모터는 어깨, 팔꿈치, 손목 같은 관절을 돌리거나 밀어내는 명령을 받는다.

따라서 로봇 제어에서는 늘 두 언어 사이를 오간다.
하나는 관절이 몇 도 돌아갔는지를 말하는 관절 언어이고, 다른 하나는 손끝이 공간에서 어디에 있는지를 말하는 작업 언어이다.
이 장의 목적은 로봇공학 전체를 다루는 것이 아니라, 뒤의 임피던스 제어 장을 읽기 위한 최소한의 번역 규칙을 준비하는 것이다.
관절공간, 작업공간, 정기구학, 역기구학, 자코비안, 궤적이라는 말이 왜 필요한지부터 차례로 보자.
이 장의 기본 흐름은 Lynch와 Park의 Modern Robotics가 다루는 기구학, 속도기구학, 정역학, 궤적 생성의 표준 틀을 교육용으로 풀어 쓴 것이다 \cite{lynchpark2017modernrobotics}.

\paragraph{이 장을 읽는 순서.}
\vmark{map-section-reading-order}%
처음 읽을 때는 이 장을 실제 제어기가 반드시 따르는 계산 파이프라인이 아니라, 학습 내비게이션으로 보면 된다.
\vmark{map-learning-navigation}%
\vmark{map-not-pipeline}%
먼저 \(\bm{q}\)라는 관절 언어를 정하고, 정기구학 \(\bm{x}=f(\bm{q})\)으로 관절값에서 손끝 위치를 읽는다.
\vmark{map-joint-language-fk}%
역기구학은 원하는 손끝 위치에서 가능한 관절값 후보를 찾는 반대 질문이다.
그다음 자코비안은 현재 자세 근처에서 작은 변화와 속도를 번역한다.
\vmark{map-jacobian-small-change}%
즉 \(\delta\bm{x}\approx\bm{J}\delta\bm{q}\), \(\dot{\bm{x}}=\bm{J}\dot{\bm{q}}\)를 통해 관절 쪽 변화가 손끝 쪽 변화로 어떻게 보이는지 읽는다.
감쇠 최소제곱(DLS)은 특이점 근처에서 원하는 손끝 속도를 관절 속도로 되돌릴 때 관절속도 폭주를 줄이는 완화식이다.
\vmark{map-dls-expanded-first-use}%
\vmark{map-dls-relief}%
\(\bm{J}^{\mathsf{T}}\bm{f}\)는 속도 역계산이 아니라 손끝 힘을 관절 effort로 읽는 힘 번역이다.
\vmark{map-jtf-force-translation}%
마지막으로 경로와 궤적은 관절공간 또는 작업공간 명령에 시간표를 붙여 속도, 가속도, 저크 요구를 읽는 관점이다.
\vmark{map-path-trajectory-timetable}%

\subsection{링크, 관절, 관절 변수}
\label{sec:robotics-links-joints-variables}

로봇 팔은 대체로 단단한 막대나 몸체가 관절로 이어진 구조이다.
단단한 몸체를 링크(link)라고 부르고, 링크와 링크 사이에서 허용된 움직임을 관절(joint)이라고 부른다.
문짝과 경첩을 떠올리면 가장 쉽다.
문짝은 링크이고, 경첩은 문짝이 한 축을 중심으로 회전하도록 허용하는 관절이다.

산업용 매니퓰레이터에서 가장 자주 만나는 관절은 두 가지이다.
\begin{itemize}
  \item 회전 관절(revolute joint): 경첩처럼 한 축을 중심으로 돈다. 관절 변수는 각도이고 보통 라디안(rad)으로 쓴다.
  \item 직선 관절(prismatic joint): 서랍 레일처럼 한 방향으로 밀려 나가거나 들어온다. 관절 변수는 길이이고 보통 미터(m)로 쓴다.
\end{itemize}

각 관절의 값을 모아 쓴 벡터가 관절 변수이다.
관절이 \(n\)개인 로봇이면
\begin{equation}
  \bm{q}
  =
  \begin{bmatrix}
    q_1 & q_2 & \cdots & q_n
  \end{bmatrix}^{\mathsf{T}}
  \in \R^n
\end{equation}
처럼 쓴다.
여기서 \(q_i\)는 \(i\)번째 관절의 각도 또는 직선 변위이다.
회전 관절이면 \(q_i\)의 단위가 rad이고, 직선 관절이면 m이다.
전기 회로에서 전하를 \(q\)로 썼던 것과 같은 글자이지만, 로봇 장에서는 굵은 \(\bm{q}\) 또는 아래첨자 \(q_i\)를 관절 변수로 읽는다.

관절 속도, 관절 가속도, 관절 저크(jerk)는 시간 미분으로 정의한다.
\begin{equation}
  \dot{\bm{q}}=\frac{\dd \bm{q}}{\dd t},
  \qquad
  \ddot{\bm{q}}=\frac{\dd^2 \bm{q}}{\dd t^2},
  \qquad
  \dddot{\bm{q}}=\frac{\dd^3 \bm{q}}{\dd t^3}.
\end{equation}
\(\dot{\bm{q}}\)는 각속도 또는 직선 속도, \(\ddot{\bm{q}}\)는 각가속도 또는 직선 가속도, \(\dddot{\bm{q}}\)는 가속도가 얼마나 갑자기 바뀌는지를 나타내는 값이다.
운전할 때 속도 자체보다 갑자기 가속 페달을 밟거나 떼는 순간이 몸에 더 거칠게 느껴지는 것처럼, 로봇에서도 저크가 큰 명령은 진동, 소음, 구동기 부담을 키울 수 있다.

\begin{table}[!htbp]
  \centering
  \caption{관절 변수와 시간 미분의 물리적 의미}
  \label{tab:joint-variable-derivatives}
  \small
  \renewcommand{\arraystretch}{1.12}
  \begin{tabularx}{\linewidth}{@{}p{0.18\linewidth}p{0.25\linewidth}X@{}}
    \toprule
    기호 & 이름 & 쉽게 읽는 법 \\
    \midrule
    \(\bm{q}\) & 관절 위치 & 각 관절이 얼마나 돌았거나 밀렸는가 \\
    \(\dot{\bm{q}}\) & 관절 속도 & 각 관절이 지금 얼마나 빠르게 움직이는가 \\
    \(\ddot{\bm{q}}\) & 관절 가속도 & 관절 속도가 얼마나 빠르게 변하는가 \\
    \(\dddot{\bm{q}}\) & 관절 저크 & 가속도가 얼마나 갑자기 바뀌는가 \\
    \(\bm{\tau}\) & 관절 힘/토크 & 회전 관절에서는 토크, 직선 관절에서는 힘 \\
    \bottomrule
  \end{tabularx}
\end{table}

여기서 configuration과 state를 한 번 구분해 두면 뒤의 제어식이 덜 헷갈린다.
Configuration은 기하학적 자세를 정하는 좌표이고, 이 장의 고정 베이스 직렬 매니퓰레이터에서는 보통 \(\bm{q}\)가 configuration이다.
\vmark{cspace-fixed-base-serial}%
\vmark{state-configuration-pose-coordinate}%
\vmark{state-fixed-base-q-configuration}%
반면 제어와 동역학에서 상태(state)는 보통 \(\bm{q}\)와 \(\dot{\bm{q}}\)를 함께 둔다.
\vmark{state-q-qdot}%
가장 흔한 강체 2차 모델에서는 미래 운동은 위치와 속도에 의존하지만, 실제 예측에는 입력, 모델, 제약조건도 함께 필요하다.
\vmark{state-future-motion}%
\(\ddot{\bm{q}}\)는 동역학과 제어 입력으로 정해지는 값으로 읽고, 최소 first-order state에 보통 따로 넣지 않는다.
\vmark{state-minimal-first-order}%
\vmark{state-qddot-input-dynamics}%
물론 구동기 동역학, 필터, 접촉 모드, 유연체, 추정기까지 모델에 넣으면 상태를 더 크게 잡을 수 있다.
\vmark{state-augmented-caveat}%
마찬가지로 저크는 궤적의 부드러움 요구나 명령의 급격함을 보는 값이지, 로봇의 기하학적 자세를 정하는 새 좌표는 아니다.
\vmark{state-jerk-smoothness}%
\vmark{state-not-geometric-coordinate}%

\(\bm{\tau}\)를 관절 토크라고만 부르면 회전 관절에서는 자연스럽지만, 직선 관절에서는 살짝 어색해진다.
정확히는 \(\tau_i\)를 \(q_i\)와 짝을 이루는 일반화 힘(generalized effort)으로 읽는 편이 안전하다.
회전 관절의 \(q_i\)는 각도이므로 그 관절의 effort는 토크이고, 직선 관절의 \(q_i\)는 길이이므로 그 관절의 effort는 선형 힘이다.
이 구분은 단위 암기가 아니라 일률을 맞추기 위한 약속에서 나온다.
\[
  P_i=\tau_i\dot{q}_i
\]
가 와트(W), 즉 초당 일의 양이 되려면 \(q_i\)가 무엇인지에 따라 \(\tau_i\)의 단위도 함께 바뀌어야 한다.
라디안은 SI 단위계에서 무차원처럼 처리되지만, 로봇 제어를 배울 때는 각도 단위로 남겨 두고 읽는 편이 실수를 줄인다.

\begin{table}[!htbp]
  \centering
  \caption{관절 종류에 따라 달라지는 관절 effort의 단위}
  \label{tab:joint-type-units-generalized-effort}
  \small
  \renewcommand{\arraystretch}{1.12}
  \begin{tabularx}{\linewidth}{@{}p{0.19\linewidth}p{0.18\linewidth}p{0.19\linewidth}X@{}}
    \toprule
    관절 종류 & \(q_i\)의 단위 & \(\tau_i\)의 단위 & 일률로 확인하는 법 \\
    \midrule
    회전 관절 & \(\mathrm{rad}\) & \(\mathrm{N\,m}\) & \(P_i=\tau_i\dot{q}_i\). 모터가 축을 돌리므로 effort는 토크이다 \\
    직선 관절 & \(\mathrm{m}\) & \(\mathrm{N}\) & \(P_i=\tau_i\dot{q}_i\). 축을 따라 밀거나 당기므로 effort는 힘이다 \\
    \bottomrule
  \end{tabularx}
\end{table}

자유도(degree of freedom, DoF)는 독립적으로 정할 수 있는 움직임의 개수이다.
단순한 회전 관절 하나는 보통 1자유도이고, 7자유도 Panda 로봇은 일곱 개의 독립 관절 변수를 갖는다.
자유도가 많다는 말은 손끝을 움직일 수 있는 선택지가 많다는 뜻이지만, 동시에 같은 손끝 위치를 만드는 관절 자세도 여러 개일 수 있다는 뜻이다.
이 점이 역기구학과 null-space, 특이점 이야기를 어렵게 만든다.

\subsection{관절공간, 작업공간, 작업영역}
\label{sec:robotics-joint-task-workspace}

관절공간(joint space)은 로봇을 관절 변수 \(\bm{q}\)로 보는 공간이다.
2자유도 평면 팔이라면 \(\bm{q}=[q_1\ q_2]^{\mathsf{T}}\)이고, 이 공간의 한 점은 어깨와 팔꿈치 각도 한 쌍을 뜻한다.
7자유도 로봇이라면 관절공간은 7차원이다.
우리가 머릿속으로 7차원 공간을 그리기는 어렵지만, 의미는 단순하다.
각 축이 관절 하나의 값이다.

여기서 Modern Robotics에서 자주 나오는 configuration space, 또는 C-space라는 말도 함께 연결해 두자.
Configuration은 로봇의 기하학적 상태를 정하는 데 필요한 좌표 묶음이다.
이 논문에서 주로 다루는 고정 베이스 직렬 매니퓰레이터에서는 관절 변수 \(\bm{q}\)가 보통 configuration을 이루므로, 관절공간을 그 로봇 팔의 configuration space로 읽어도 된다.
\vmark{cspace-joint-space-scope}%
다만 이것은 언제나 configuration space와 joint space가 같은 말이라는 뜻은 아니다.
\vmark{cspace-not-universal-identity}%
이동 베이스, 부유 베이스, 물체 포즈, 폐사슬 또는 접촉 제약까지 모델에 넣으면 전체 configuration space는 base/object pose 또는 constraints를 포함할 수 있다.
\vmark{cspace-base-object-constraints}%
반대로 task space는 우리가 선택한 출력 변수의 공간이고, workspace는 그중 로봇이 실제로 도달 가능한 작업 위치들의 집합이다.
\vmark{cspace-task-space-output}%
\vmark{cspace-workspace-reachable}%
여기서는 먼저 \(\bm{q}\)-공간과 \(\bm{x}\)-공간의 번역에 집중한다.
\vmark{cspace-q-x-translation}%

작업공간(task space)은 사람이 관심 있는 작업 변수를 모은 공간이다.
이 글처럼 손끝의 병진 위치를 \([x\ y\ z]^{\mathsf{T}}\)로 고르면, 그 작업공간은 직교 위치 공간(Cartesian position space)이 된다.
예를 들어 손끝 위치만 보면
\begin{equation}
  \bm{x}
  =
  \begin{bmatrix}
    x & y & z
  \end{bmatrix}^{\mathsf{T}}
  \in \R^3
\end{equation}
이다.
손끝의 방향까지 포함하면 위치뿐 아니라 자세도 필요하다.
자세는 보통 회전행렬 \(R\in SO(3)\) 또는 포즈 \(T\in SE(3)\)로 표현하고, 롤--피치--요처럼 3개 숫자로 쓰는 방식은 편리하지만 국소 좌표라 특이점이 생길 수 있다.
속도를 다룰 때는 트위스트(twist), 힘과 모멘트를 함께 다룰 때는 렌치(wrench)라는 표현도 등장한다.
지금은 이름만 기억해도 충분하다.
이 논문에서는 임피던스 직관을 먼저 잡기 위해 주로 손끝의 병진 위치 \(\bm{x}\)에 집중한다.

작업영역(workspace)은 선택한 작업공간 안에서 로봇 손끝이 실제로 도달할 수 있는 위치들의 집합이다.
관절공간과 작업공간은 비슷해 보이지만 다르다.
관절공간의 직선 경로는 관절각을 일정한 비율로 바꾸는 길이고, 작업공간의 직선 경로는 손끝이 실제 공간에서 직선으로 움직이는 길이다.
관절공간에서 직선으로 움직인다고 해서 손끝도 직선으로 움직이는 것은 아니다.
이 차이를 모르면 로봇이 왜 예상과 다른 곡선으로 움직이는지 이해하기 어렵다.

\begin{table}[!htbp]
  \centering
  \caption{관절공간과 작업공간의 차이}
  \label{tab:joint-task-space}
  \small
  \renewcommand{\arraystretch}{1.12}
  \begin{tabularx}{\linewidth}{@{}p{0.18\linewidth}p{0.31\linewidth}X@{}}
    \toprule
    관점 & 좌표 & 질문 \\
    \midrule
    관절공간 & \(\bm{q},\dot{\bm{q}},\ddot{\bm{q}}\) & 각 모터가 얼마만큼 움직여야 하는가? \\
    작업공간 & \(\bm{x},\dot{\bm{x}},\ddot{\bm{x}}\) 또는 포즈/트위스트 & 손끝이 공간에서 어디로, 얼마나 빠르게 움직이는가? \\
    작업영역 & 도달 가능한 \(\bm{x}\)의 집합 & 이 손끝 목표가 로봇이 실제로 갈 수 있는 곳인가? \\
    \bottomrule
  \end{tabularx}
\end{table}

접촉 작업에서는 작업공간 표현이 특히 중요하다.
벽을 누르는 방향, 표면을 따라 미끄러지는 방향, 삽입 방향처럼 사람이 말하는 작업은 대부분 손끝 기준 방향으로 표현되기 때문이다.
하지만 실제 명령은 관절 모터로 들어간다.
뒤에서 볼 자코비안은 바로 이 두 언어 사이의 순간 번역기이다.

\subsection{정기구학: 관절값에서 손끝 위치로}
\label{sec:robotics-forward-kinematics}

정기구학(forward kinematics)은 관절값 \(\bm{q}\)가 주어졌을 때 손끝 위치나 포즈를 계산하는 일이다.
가장 짧게 쓰면
\begin{equation}
  \bm{x}=f(\bm{q})
  \label{eq:fk-minimal}
\end{equation}
이다.
여기서 \(f\)는 로봇의 링크 길이, 관절 배치, 좌표계 정의를 모두 담고 있는 함수이다.

2자유도 평면 팔을 예로 보자.
첫 번째 링크 길이를 \(l_1\), 두 번째 링크 길이를 \(l_2\), 관절각을 \(q_1,q_2\)라고 하자.
어깨 관절에서 첫 링크가 \(q_1\)만큼 돌고, 둘째 링크는 첫 링크에 대해 \(q_2\)만큼 더 돈다.
그러면 손끝 위치는
\begin{align}
  x(q_1,q_2)
  &=
  l_1\cos q_1
  +
  l_2\cos(q_1+q_2),
  \label{eq:two-link-fk-x}
  \\
  y(q_1,q_2)
  &=
  l_1\sin q_1
  +
  l_2\sin(q_1+q_2).
  \label{eq:two-link-fk-y}
\end{align}
이 된다.
식이 길어 보여도 의미는 간단하다.
첫 링크 끝점의 \(x,y\) 좌표에 둘째 링크가 더해지는 것이다.
코사인은 가로 성분, 사인은 세로 성분을 만든다고 생각하면 된다.

정기구학은 보통 한 방향으로는 잘 정의된다.
관절각을 정확히 하나 정하면, 이상적인 강체 모델에서 손끝 위치도 하나로 정해진다.
그래서 시뮬레이터나 로봇 제어기 안에서는 현재 관절각을 읽어 손끝 위치를 계산하는 데 계속 사용된다.

\subsection{역기구학: 손끝 목표에서 관절값으로}
\label{sec:robotics-inverse-kinematics}

역기구학(inverse kinematics, IK)은 반대 질문이다.
원하는 손끝 위치 \(\bm{x}_d\)가 있을 때 어떤 관절값 \(\bm{q}\)가 필요한가?
기호만 보면
\begin{equation}
  \bm{q}=f^{-1}(\bm{x}_d)
\end{equation}
처럼 쓰고 싶지만, 실제로는 조심해야 한다.
정기구학 함수 \(f\)는 보통 단순한 일대일 함수가 아니기 때문이다.

역기구학에서는 세 가지 일이 생길 수 있다.
\begin{itemize}
  \item 해가 없다. 목표점이 작업영역 밖에 있으면 어떤 관절각으로도 갈 수 없다.
  \item 해가 하나 있다. 특별한 조건에서 관절 자세가 하나로 정해질 수 있다.
  \item 해가 여러 개 있다. 2링크 팔에서도 같은 손끝 위치에 대해 팔꿈치를 위로 접는 해와 아래로 접는 해가 생길 수 있다.
\end{itemize}

2링크 평면 팔에서는 목표점까지의 거리 \(r\)만 보아도 먼저 해의 개수를 대략 판정할 수 있다.
\begin{table}[!htbp]
  \centering
  \caption{2링크 위치 IK에서 거리로 먼저 보는 해의 개수}
  \label{tab:two-link-ik-reachability}
  \begin{tabularx}{\linewidth}{>{\raggedright\arraybackslash}p{0.28\linewidth}>{\raggedright\arraybackslash}X>{\raggedright\arraybackslash}p{0.18\linewidth}}
    \toprule
    조건 & 의미 & IK 해석 \\
    \midrule
    \(r>l_1+l_2\) & 팔을 모두 펴도 닿지 않음 & 해 없음 \\
    \(r<|l_1-l_2|\) & 한 링크가 너무 길어 안쪽 빈 구멍에 들어감 & 해 없음 \\
    \(r=l_1+l_2\) 또는 \(r=|l_1-l_2|\) & 완전히 펴지거나 접힌 경계 & 경계 해 하나 \\
    \(|l_1-l_2|<r<l_1+l_2\) & 두 링크가 삼각형을 만들 수 있음 & 보통 두 branch \\
    \bottomrule
  \end{tabularx}
\end{table}

2링크 팔에서는 이 세 경우가 어디서 나오는지 코사인 법칙으로 한 번 볼 수 있다.
목표 손끝 위치를 \((x_d,y_d)\)라고 하고, 어깨에서 목표점까지의 거리 제곱을
\begin{equation}
  r^2=x_d^2+y_d^2
\end{equation}
라고 두자.
어깨, 팔꿈치, 손끝이 만드는 삼각형에서 두 링크 길이 \(l_1,l_2\)와 거리 \(r\)가 정해지면, 팔꿈치 각도 \(q_2\)는
\begin{equation}
  \cos q_2
  =
  \frac{r^2-l_1^2-l_2^2}{2l_1l_2}
  \label{eq:two-link-ik-cos-q2}
\end{equation}
를 만족해야 한다.
오른쪽 값이 \(1\)보다 크거나 \(-1\)보다 작으면 실제 각도 \(q_2\)가 존재하지 않으므로 해가 없다.
오른쪽 값이 정확히 \(1\)이나 \(-1\)이면 팔이 완전히 펴지거나 접힌 경계 자세라 해가 하나로 붙는다.
다만 여기서의 하나는 보통의 비퇴화 2링크 기하와 관절 한계를 잠시 무시한 말이다.
예를 들어 \(l_1=l_2\)이고 목표점이 어깨 원점에 놓이는 특수한 경우에는 \(q_2=\pi\)로 접힌 채 \(q_1\)을 여러 방향으로 돌릴 수 있고, 반대로 실제 로봇의 관절 한계가 있으면 수식상 가능한 branch가 사라질 수도 있다.
그 사이 값이면 \(q_2=\arccos(\cdot)\)와 \(q_2=-\arccos(\cdot)\)가 모두 가능해서 팔꿈치 위쪽 해와 아래쪽 해가 생긴다.
\(\arccos\) 자체는 \(0\)부터 \(\pi\) 사이의 대표 각도 하나를 돌려주므로, 같은 cosine 값을 갖는 두 자세를 표현하려면 부호를 따로 붙여 읽어야 한다.
이 작은 계산만 보아도 IK가 단순히 \(f^{-1}\) 하나로 끝나는 문제가 아니라는 점이 드러난다.

실제로 계산할 때는 다음 순서로 읽으면 된다.
먼저
\begin{equation}
  c
  =
  \frac{r^2-l_1^2-l_2^2}{2l_1l_2}
  \label{eq:two-link-ik-c-value}
\end{equation}
를 계산한다.
\(c\)가 \([-1,1]\) 범위 안에 있으면 팔꿈치 각도 후보는
\begin{equation}
  q_2^{(+)}
  =
  \arccos c,
  \qquad
  q_2^{(-)}
  =
  -\arccos c
  \label{eq:two-link-ik-q2-branches}
\end{equation}
처럼 두 개가 된다.
그 다음에는 목표점 방향각에서 팔 안쪽 삼각형의 보정각을 빼서 \(q_1\)을 구한다.
\begin{equation}
  q_1^{(\pm)}
  =
  \operatorname{atan2}(y_d,x_d)
  -
  \operatorname{atan2}
  \left(
    l_2\sin q_2^{(\pm)},
    l_1+l_2\cos q_2^{(\pm)}
  \right).
  \label{eq:two-link-ik-q1-from-q2}
\end{equation}
여기서 \(\operatorname{atan2}(y,x)\)는 점 \((x,y)\)가 원점에서 보이는 방향각을 사분면까지 고려해서 돌려주는 함수이다.
보통 \(\tan^{-1}(y/x)\)만 쓰면 \(x\)가 음수인 경우 사분면을 헷갈릴 수 있기 때문에, 로봇공학 코드에서는 \(\operatorname{atan2}\)를 쓰는 편이 안전하다.

숫자를 하나 넣어 보자.
\(l_1=l_2=1~\mathrm{m}\), 목표점 \((x_d,y_d)=(1,1)~\mathrm{m}\)이면 \(r^2=2\)이고 \(c=0\)이다.
따라서
\begin{align}
  q_2^{(+)}
  &=
  \frac{\pi}{2},
  &
  q_1^{(+)}
  &=
  \operatorname{atan2}(1,1)
  -
  \operatorname{atan2}(1,1)
  =
  0,
  \label{eq:two-link-ik-positive-branch}
  \\
  q_2^{(-)}
  &=
  -\frac{\pi}{2},
  &
  q_1^{(-)}
  &=
  \operatorname{atan2}(1,1)
  -
  \operatorname{atan2}(-1,1)
  =
  \frac{\pi}{2}.
  \label{eq:two-link-ik-negative-branch}
\end{align}
두 해를 다시 정기구학 식에 넣으면 둘 다 같은 손끝점 \((1,1)\)을 만든다.
다만 첫 번째 해 \((q_1,q_2)=(0,\pi/2)\)에서는 팔꿈치가 \((1,0)\)에 있고, 두 번째 해 \((q_1,q_2)=(\pi/2,-\pi/2)\)에서는 팔꿈치가 \((0,1)\)에 있다.
손끝은 같지만 팔꿈치가 놓인 쪽이 다른 것이다.
그림~\ref{fig:two-link-ik-branches}처럼 두 branch를 같은 좌표계에 겹쳐 그리면, 같은 목표점에 도달하더라도 팔꿈치가 어깨--목표선의 서로 다른 쪽에 놓인다는 점이 보인다.

\begin{figure}[!htbp]
  \centering
  \resizebox{\linewidth}{!}{%
  \begin{tikzpicture}[
      axis/.style={-{Latex[length=1.7mm]}, draw=black!65, line width=0.45pt},
      guide/.style={densely dashed, draw=black!35, line width=0.45pt},
      panel/.style={draw=black!25, rounded corners=2pt, fill=black!2},
      poslink/.style={draw=blue!70!black, line width=1.25pt},
      neglink/.style={draw=orange!85!black, line width=1.25pt, dashed},
      point/.style={circle, fill=black, inner sep=1.5pt},
      target/.style={circle, fill=green!45!black, inner sep=1.8pt},
      lab/.style={font=\scriptsize, align=center},
      title/.style={font=\scriptsize\bfseries, align=center},
      note/.style={draw=black!30, rounded corners=2pt, fill=white, align=left, font=\scriptsize, inner sep=4pt}
    ]

    \draw[panel] (-0.35,-0.35) rectangle (8.9,3.45);

    \begin{scope}[shift={(0,0)}]
      \draw[axis] (-0.08,0) -- (2.25,0) node[right, lab] {$x$};
      \draw[axis] (0,-0.08) -- (0,2.25) node[above, lab] {$y$};

      \coordinate (S) at (0,0);
      \coordinate (Ep) at (1,0);
      \coordinate (En) at (0,1);
      \coordinate (T) at (1,1);

      \draw[guide] (S) -- (T);
      \node[lab, rotate=45, anchor=south] at (0.54,0.54) {기준선 \(\bm{x}_d\)};

      \draw[poslink] (S) -- (Ep) -- (T);
      \draw[neglink] (S) -- (En) -- (T);
      \node[lab, text=blue!70!black, below] at (0.5,-0.05) {$l_1$};
      \node[lab, text=blue!70!black, right] at (1.05,0.5) {$l_2$};
      \node[lab, text=orange!85!black, left] at (-0.05,0.5) {$l_1$};
      \node[lab, text=orange!85!black, above] at (0.5,1.05) {$l_2$};

      \node[point, label={[lab, below left]$S=(0,0)$}] at (S) {};
      \node[target, label={[lab, above right]$T=(1,1)$}] at (T) {};
      \node[point, blue!70!black, fill=blue!70!black,
        label={[lab, below]$E_+=(1,0)$}] at (Ep) {};
      \node[point, orange!85!black, fill=orange!85!black,
        label={[lab, left]$E_-=(0,1)$}] at (En) {};

    \end{scope}

    \begin{scope}[shift={(3.1,0.12)}]
      \node[title, anchor=west] at (0,2.95) {같은 손끝점, 다른 branch};

      \node[note, text width=5.35cm, anchor=north west] at (0,2.72)
        {\textbf{손끝은 같다}\\
        두 해 모두 정기구학에 넣으면 \(T=(1,1)\)을 만든다.\\[2pt]
        \textbf{팔꿈치 위치는 다르다}\\
        \(q_2>0\): \(E_+=(1,0)\), \(s_+=x_d e_y-y_d e_x=-1\).\\
        \(q_2<0\): \(E_-=(0,1)\), \(s_-=x_d e_y-y_d e_x=+1\).\\[2pt]
        \textbf{주의}\\
        위/아래 이름은 관례이므로 branch 정의를 명시한다.};
    \end{scope}
  \end{tikzpicture}%
  }
  \caption{관절 한계와 \(2\pi\) 각도 wrapping을 잠시 무시한 2R 팔에서 같은 목표점 \((1,1)\)에 도달하는 두 대표 IK branch이다. 여기서는 \(s=x_d e_y-y_d e_x\)로 팔꿈치가 어깨--목표 기준선의 어느 쪽에 있는지 정하므로, \(q_2>0\) branch는 \(s<0\), \(q_2<0\) branch는 \(s>0\)에 놓인다. 따라서 팔꿈치 위쪽/아래쪽이라는 말은 좌표계와 그림 관례에 의존하며, 실험이나 코드에서는 \(q_2\)의 부호 또는 signed side로 branch를 분명히 정의하는 편이 안전하다.}
  \label{fig:two-link-ik-branches}
\end{figure}


어깨에서 목표점으로 향하는 선을 기준으로 팔꿈치가 어느 쪽에 있는지 숫자로 보려면, 팔꿈치 좌표를 \((e_x,e_y)\)라고 두고 \(x_d e_y-y_d e_x\)의 부호를 볼 수 있다.
이 예제에서는 \(q_2>0\) branch가 \(-1\), \(q_2<0\) branch가 \(+1\)이므로 서로 반대쪽에 있다.
그래서 팔꿈치 위쪽/아래쪽이라는 이름은 그림의 좌표계와 관례에 따라 달라질 수 있고, 코드나 실험 설정에서는 \(q_2\)의 부호 또는 upper/lower branch 정의를 명시하는 것이 안전하다.
여기서 branch는 같은 목표점에 도달하는 연속적인 자세 선택의 갈래를 뜻한다.
두 해는 손끝은 같지만 팔꿈치 위치가 다르므로, 이후 같은 손끝 명령도 서로 다른 관절 속도와 effort로 번역될 수 있다.

여기서 해가 여러 개 있다는 말은 단순히 수학 풀이가 귀찮다는 뜻이 아니다.
같은 손끝 위치라도 로봇의 팔꿈치, 손목, 어깨 자세가 달라질 수 있고, 그 결과 관절 속도, 관절 토크, 특이점까지의 거리도 달라진다.
2자유도 평면 팔에서는 이것이 팔꿈치 위쪽 해와 아래쪽 해처럼 보인다.
손끝은 같은 점에 있어도 한 자세는 팔꿈치가 위로 접혀 있고, 다른 자세는 아래로 접혀 있을 수 있다.
Lab03의 upper/lower 또는 mirrored branch 비교는 바로 이 차이를 플롯으로 읽기 위한 장면이다.
목표 손끝 경로가 비슷해도 관절각 곡선, 조건수, effort 지표가 다르게 나타날 수 있다.
따라서 IK의 어려움은 목표에 도달하느냐만의 문제가 아니라, 여러 해 중 어떤 자세 가지(branch)를 선택하느냐의 문제이기도 하다.

자유도가 작업 변수보다 많은 로봇에서는 이 문제가 더 자연스럽게 커진다.
예를 들어 손끝 병진 위치만 \((x,y,z)\) 세 숫자로 정하는데 로봇에는 관절이 7개 있다면, 남는 선택지가 생긴다.
이 남는 관절 움직임을 보통 여유 자유도(redundancy)라고 부르고, 손끝 위치를 거의 바꾸지 않는 방향의 관절 움직임을 null-space 움직임이라고 부른다.
2링크 위치 IK에서는 보통 두 개의 떨어진 branch를 고르는 문제로 보이지만, 7자유도 로봇에서는 해가 연속적인 가족처럼 생길 수 있다.
직관적으로 말하면 손끝은 같은 곳에 두고 팔꿈치 모양만 조금 바꾸는 여지가 있는 셈이다.
이 여유는 관절 한계를 피하거나, 특이점에서 멀어지거나, 더 편한 자세를 고르는 데 사용할 수 있다.
하지만 이 논문과 Lab04에서는 null-space 최적화 자체를 깊게 다루기보다, 7자유도 로봇에서는 같은 손끝 목표라도 현재 자세와 역기구학 선택에 따라 관절공간의 부담이 달라질 수 있다는 점을 기억하는 정도로 둔다.
즉, 여기서의 핵심은 더 많은 관절이 항상 더 쉽다는 말이 아니라, 선택지가 늘어난 만큼 어떤 관절 자세를 고를 것인지도 제어 문제의 일부가 된다는 점이다.

이 점은 임피던스 제어에서도 중요하다.
손끝에 같은 가상 스프링을 걸어도, 어떤 관절 자세로 그 손끝 위치를 만들고 있는지에 따라 필요한 관절 움직임과 effort가 달라진다.
그래서 Lab03의 2자유도 팔은 단순한 장난감 예제가 아니라, 관절공간과 작업공간 사이의 번역 문제가 관절각, 속도, effort, 조건수에 어떤 흔적을 남기는지 보여주는 중간 계단이다.

\paragraph{IK에서 자코비안으로 넘어가는 다리.}
\vmark{handoff-ik-jacobian-bridge}%
가지 선택 뒤의 읽기 순서를 한 번만 정리하자.
\vmark{handoff-branch-reading-order}%
IK는 가능한 branch 후보를 보여 주지만, 이후 계산은 실제로 선택된 현재 branch의 \(\bm{q}\)에서 시작한다.
\vmark{handoff-current-branch-q}%
그 \(\bm{q}\)에서 \(\bm{J}(\bm{q})\)를 계산하므로, 같은 손끝 위치라도 branch가 다르면 자코비안, 조건수, 관절속도, effort가 달라질 수 있다.
\vmark{handoff-different-branch-metrics}%
\vmark{handoff-jacobian-at-that-q}%
DLS는 branch를 고르는 마법이 아니라, 선택된 현재 자세 근처에서 원하는 손끝 속도를 너무 큰 관절속도 없이 만들기 위한 완화식이다.
\vmark{handoff-dls-not-branch-magic}%
마찬가지로 \(\bm{J}^{\mathsf{T}}\bm{f}\)도 같은 현재 자세에서 힘을 관절 effort로 읽는다.
\vmark{handoff-same-pose-effort}%
그래서 다음 절의 자코비안은 손끝 목표와 선택된 현재 branch의 \(\bm{q}\)를 정한 뒤, 그 자세 근처에서 속도와 힘이 어떻게 번역되는지를 묻는 단계라고 보면 된다.
\vmark{handoff-branch-pose-framing}%
\vmark{handoff-velocity-force-translation}%

\subsection{자코비안 1단계: 작은 변화의 장부}
\label{sec:robotics-jacobian-small-motion}

자코비안(Jacobian)은 처음 보면 어려운 이름처럼 들리지만, 출발점은 아주 단순하다.
정기구학이
\[
  \bm{x}=f(\bm{q})
\]
라면, 자코비안은 관절을 아주 조금 움직였을 때 손끝이 얼마나 조금 움직이는지를 모은 표이다.

2자유도 평면 팔에서 손끝 좌표가 \(x(q_1,q_2),y(q_1,q_2)\)라고 하자.
관절각이 아주 조금 \(\Delta q_1,\Delta q_2\)만큼 바뀌면 손끝 위치도 조금 바뀐다.
그 작은 변화는 1차 근사로
\begin{align}
  \Delta x
  &\approx
  \frac{\partial x}{\partial q_1}\Delta q_1
  +
  \frac{\partial x}{\partial q_2}\Delta q_2,
  \label{eq:jacobian-dx}
  \\
  \Delta y
  &\approx
  \frac{\partial y}{\partial q_1}\Delta q_1
  +
  \frac{\partial y}{\partial q_2}\Delta q_2.
  \label{eq:jacobian-dy}
\end{align}
여기서 \(\partial x/\partial q_1\)은 \(q_1\)만 아주 조금 바꿨을 때 \(x\)가 얼마나 민감하게 변하는지를 뜻한다.
마찬가지로 \(\partial y/\partial q_2\)는 \(q_2\)만 바꿨을 때 \(y\)가 얼마나 변하는지를 뜻한다.

이 네 개의 민감도를 행렬로 모으면
\begin{equation}
  \bm{J}(\bm{q})
  =
  \begin{bmatrix}
    \partial x/\partial q_1 & \partial x/\partial q_2 \\
    \partial y/\partial q_1 & \partial y/\partial q_2
  \end{bmatrix}
  \label{eq:jacobian-partial-table}
\end{equation}
가 된다.
그러면 식 \eqref{eq:jacobian-dx}, \eqref{eq:jacobian-dy}를 한 줄로
\begin{equation}
  \Delta \bm{x}
  \approx
  \bm{J}(\bm{q})\Delta\bm{q}
  \label{eq:jacobian-small-motion}
\end{equation}
라고 쓸 수 있다.

이 식은 자코비안을 이해하는 첫 번째 핵심이다.
자코비안은 큰 이동 전체를 정확히 푸는 마법 공식이 아니라, 현재 자세 근처에서 작은 관절 변화가 작은 손끝 변화로 어떻게 바뀌는지 알려주는 국소적인 선형 근사이다.
그래서 자코비안은 항상 현재 자세 \(\bm{q}\)에 따라 달라진다.
팔꿈치가 접혀 있을 때와 쭉 펴져 있을 때, 같은 관절 변화가 손끝에 만드는 움직임은 다르다.

2링크 팔의 정기구학 식 \eqref{eq:two-link-fk-x}, \eqref{eq:two-link-fk-y}에서 직접 편미분해 보자.
먼저 \(x\)식은 코사인으로 되어 있으므로, 미분하면 \(-\sin\)이 나온다.
\(q_1+q_2\)를 \(q_1\)로 미분해도 1이고, \(q_2\)로 미분해도 1이지만, 첫 번째 링크 항 \(l_1\cos q_1\)은 \(q_2\)를 포함하지 않으므로 \(q_2\)에 대해서는 사라진다.
\begin{align}
  \frac{\partial x}{\partial q_1}
  &=
  -l_1\sin q_1-l_2\sin(q_1+q_2),
  &
  \frac{\partial x}{\partial q_2}
  &=
  -l_2\sin(q_1+q_2),
  \\
  \frac{\partial y}{\partial q_1}
  &=
  l_1\cos q_1+l_2\cos(q_1+q_2),
  &
  \frac{\partial y}{\partial q_2}
  &=
  l_2\cos(q_1+q_2).
\end{align}
이 네 항을 표처럼 모으면
\begin{equation}
  \bm{J}(\bm{q})
  =
  \begin{bmatrix}
    -l_1\sin q_1-l_2\sin(q_1+q_2)
    &
    -l_2\sin(q_1+q_2)
    \\
    l_1\cos q_1+l_2\cos(q_1+q_2)
    &
    l_2\cos(q_1+q_2)
  \end{bmatrix}.
  \label{eq:two-link-jacobian}
\end{equation}
첫 번째 열은 \(q_1\)을 단위 속도로 움직이고 \(q_2\)는 가만히 둘 때 손끝이 어느 방향으로 움직이는지를 나타낸다.
두 번째 열은 \(q_2\)를 단위 속도로 움직이고 \(q_1\)은 가만히 둘 때의 손끝 움직임이다.
즉 자코비안의 각 열은 각 관절 하나가 손끝에 기여하는 순간 움직임 벡터라고 볼 수 있다.

숫자를 하나 넣어 보면 더 선명하다.
\(l_1=l_2=1~\mathrm{m}\), \(q_1=0\), \(q_2=\pi/2\)라면 손끝은 \((1,1)\)에 있고,
\[
  \bm{J}
  =
  \begin{bmatrix}
    -1 & -1 \\
     1 &  0
  \end{bmatrix}
\]
가 된다.
그림~\ref{fig:two-link-jacobian-columns}처럼 이 두 열을 손끝에 붙은 속도 화살표로 그리면, 행렬의 각 열이 왜 특정 관절 하나의 순간 기여를 뜻하는지 바로 보인다.

\begin{figure}[!htbp]
  \centering
  \resizebox{\linewidth}{!}{%
  \begin{tikzpicture}[
      axis/.style={-{Latex[length=1.7mm]}, draw=black!65, line width=0.45pt},
      guide/.style={densely dashed, draw=black!35, line width=0.45pt},
      link/.style={draw=black!70, line width=1.2pt},
      colone/.style={-{Latex[length=2mm]}, draw=blue!70!black, line width=1.25pt},
      coltwo/.style={-{Latex[length=2mm]}, draw=orange!85!black, line width=1.25pt},
      panel/.style={draw=black!25, rounded corners=2pt, fill=black!2},
      point/.style={circle, fill=black, inner sep=1.5pt},
      hand/.style={circle, fill=green!45!black, inner sep=1.8pt},
      lab/.style={font=\scriptsize, align=center},
      title/.style={font=\scriptsize\bfseries, align=center},
      note/.style={draw=black!30, rounded corners=2pt, fill=white, align=left, font=\scriptsize, inner sep=4pt}
    ]

    \draw[panel] (-0.35,-0.35) rectangle (10.25,3.45);

    \begin{scope}[shift={(0.22,0.18)}, scale=1.18]
      \draw[axis] (-0.08,0) -- (2.35,0) node[right, lab] {$x$};
      \draw[axis] (0,-0.08) -- (0,2.35) node[above, lab] {$y$};

      \coordinate (S) at (0,0);
      \coordinate (E) at (1,0);
      \coordinate (H) at (1,1);

      \draw[link] (S) -- (E) -- (H);
      \draw[guide] (S) -- (H);
      \draw[guide] (E) -- (H);

      \draw[colone] (0.34,0) arc[start angle=0, end angle=52, radius=0.34];
      \node[lab, text=blue!70!black] at (0.53,0.40) {$+q_1$};
      \draw[coltwo] (1.25,0.20) arc[start angle=35, end angle=120, radius=0.28];
      \node[lab, text=orange!85!black] at (1.42,0.42) {$+q_2$};

      \draw[colone] (H) -- ++(-0.58,0.58)
        node[lab, text=blue!70!black, above left] {$\bm{j}_1=(-1,1)$};
      \draw[coltwo] (H) -- ++(-0.78,0)
        node[lab, text=orange!85!black, left] {$\bm{j}_2=(-1,0)$};

      \node[point, label={[lab, below left]$S$}] at (S) {};
      \node[point, label={[lab, below right]$E$}] at (E) {};
      \node[hand, label={[lab, above right]$H=(1,1)$}] at (H) {};
    \end{scope}

    \begin{scope}[shift={(3.7,0.12)}]
      \node[title, anchor=west] at (0,2.95) {자코비안 열을 손끝 속도로 읽기};

      \node[note, text width=6.05cm, anchor=north west] at (0,2.72)
        {\textbf{정규화 예제}\\
        \(l_1=l_2=1\), \(q_1=0\), \(q_2=\pi/2\).\\[2pt]
        \textbf{첫 번째 열}\\
        \(q_1\)만 \(1~\mathrm{rad/s}\)로 돌리면 손끝은 어깨 \(S\)를 중심으로 움직인다.\\
        \(S\rightarrow H=(1,1)\)의 접선 방향이 \(\bm{j}_1=(-1,1)~\mathrm{m/rad}\)이다.\\[2pt]
        \textbf{두 번째 열}\\
        \(q_2\)만 \(1~\mathrm{rad/s}\)로 돌리면 둘째 링크 끝은 팔꿈치 \(E\)를 중심으로 움직인다.\\
        \(E\rightarrow H=(0,1)\)의 접선 방향이 \(\bm{j}_2=(-1,0)~\mathrm{m/rad}\)이다.\\[2pt]
        \(\bm{J}=[\bm{j}_1\ \bm{j}_2]\).};
    \end{scope}
  \end{tikzpicture}%
  }
  \caption{정규화된 2링크 팔 자세 \(l_1=l_2=1~\mathrm{m}\), \(q_1=0\), \(q_2=\pi/2\)에서 자코비안의 두 열을 손끝 순간속도 방향으로 읽은 그림이다. 첫 번째 열 \(\bm{j}_1=(-1,1)^{\mathsf T}\)은 어깨 관절 \(q_1\)만 단위 속도로 움직이고 팔꿈치 \(q_2\)를 고정했을 때의 손끝 순간속도이다. 두 번째 열 \(\bm{j}_2=(-1,0)^{\mathsf T}\)은 팔꿈치 관절 \(q_2\)만 움직이고 어깨 \(q_1\)을 고정했을 때의 손끝 순간속도이다. 화살표는 유한한 이동 경로가 아니라 순간속도 방향을 표시하기 위해 축척해 그린 것이다. 따라서 자코비안 행렬의 열은 숫자 목록이 아니라, 현재 자세에서 각 관절 하나가 손끝에 만드는 속도 화살표로 볼 수 있다.}
  \label{fig:two-link-jacobian-columns}
\end{figure}


이제 같은 생각을 아주 작은 변화 대신 1초 동안의 변화, 즉 속도로 읽어 보자.
이때 첫 번째 관절만 \(0.1~\mathrm{rad/s}\)로 움직이면 손끝 속도는 \([-0.1\ \ 0.1]^{\mathsf{T}}~\mathrm{m/s}\)이고, 두 번째 관절만 \(0.1~\mathrm{rad/s}\)로 움직이면 \([-0.1\ \ 0]^{\mathsf{T}}~\mathrm{m/s}\)이다.
같은 \(0.1~\mathrm{rad/s}\)라도 어떤 관절을 움직이는지에 따라 손끝 방향이 달라진다는 뜻이다.

\subsection{자코비안 2단계: 속도 관계}
\label{sec:robotics-velocity-kinematics}

작은 변화 식에서 \(\Delta t\)를 점점 작게 보내는 극한을 취하면, 즉 연쇄법칙을 적용하면 속도 관계가 된다.
\begin{equation}
  \dot{\bm{x}}
  =
  \bm{J}(\bm{q})\dot{\bm{q}}.
  \label{eq:jacobian-velocity-foundation}
\end{equation}
이 식이 속도기구학(velocity kinematics)의 핵심이다.
관절 속도 \(\dot{\bm{q}}\)를 넣으면 현재 자세에서 손끝 속도 \(\dot{\bm{x}}\)가 나온다.
이때 \(\bm{J}\)의 단위도 의미가 있다.
회전 관절만 있는 평면 팔에서 \(\dot{q}\)의 단위는 \(\mathrm{rad/s}\), \(\dot{x}\)의 단위는 \(\mathrm{m/s}\)이므로, 병진 자코비안의 성분은 대략 \(\mathrm{m/rad}\)처럼 읽을 수 있다.
직선 관절이 섞이면 그 열의 단위는 \(\mathrm{m/m}\), 즉 축 방향을 나타내는 무차원 방향 벡터처럼 읽힌다.
따라서 하나의 자코비안 안에서도 회전 관절 열과 직선 관절 열은 서로 다른 단위 감각을 가질 수 있다.
이것이 뒤에서 \(\bm{J}^{\mathsf{T}}\bm{f}\)를 계산했을 때 어떤 성분은 토크가 되고, 어떤 성분은 선형 힘이 되는 이유이다.

가속도까지 가면 한 항이 더 생긴다.
여기서 \(\bm{x}\)는 먼저 고정 좌표계에서 본 손끝의 병진 위치로 읽자.
\vmark{accel-fixed-frame-translation}%
속도 식 \(\dot{\bm{x}}=\bm{J}(\bm{q})\dot{\bm{q}}\)에서 \(\bm{J}(\bm{q})\)는 고정된 숫자표가 아니다.
\vmark{accel-jacobian-not-fixed-table}%
관절각 \(\bm{q}(t)\)가 움직이면 로봇 자세가 바뀌고, 그 자세에서의 자코비안 값도 함께 바뀐다.
따라서 속도 식을 시간으로 한 번 더 미분할 때는 \(\bm{J}\)와 \(\dot{\bm{q}}\)의 곱을 미분한다고 읽어야 한다.
\vmark{accel-product-rule}%
\begin{equation}
  \ddot{\bm{x}}
  =
  \bm{J}(\bm{q})\ddot{\bm{q}}
  +
  \dot{\bm{J}}(\bm{q},\dot{\bm{q}})\dot{\bm{q}}.
  \label{eq:jacobian-acceleration-foundation}
\end{equation}
첫 항은 관절 가속도가 손끝 가속도로 바뀌는 효과이고, 둘째 항은 로봇 자세가 움직이면서 자코비안 자체도 변하기 때문에 생기는 효과이다.
짧게 쓰면
\[
  \dot{\bm{J}}
  =
  \frac{\dd}{\dd t}\bm{J}(\bm{q}(t))
\vmark{accel-jdot-definition}%
\]
라는 뜻이다.
즉 \(\dot{\bm{J}}\)는 독립적으로 새로 생긴 행렬이 아니라, 관절 속도 \(\dot{\bm{q}}\) 때문에 현재 자세가 지나가며 자코비안 표가 변하는 속도이다.
\vmark{accel-jdot-not-independent}%
회전 관절 열만 놓고 단위를 확인하면 \(\bm{J}\ddot{\bm{q}}\)는 \((\mathrm{m/rad})(\mathrm{rad/s^2})=\mathrm{m/s^2}\), \(\dot{\bm{J}}\dot{\bm{q}}\)는 \((\mathrm{m/(rad\,s)})(\mathrm{rad/s})=\mathrm{m/s^2}\)로 둘 다 손끝 가속도 단위가 된다.
\vmark{accel-unit-check}%
이 항의 크기와 방향은 현재 자세, 관절 속도, 좌표계와 양의 관절 방향 약속에 따라 달라지므로 항상 목표 쪽이나 반대쪽으로 작용하는 항으로 외우면 안 된다.
\vmark{accel-sign-direction-caveat}%
처음 읽을 때는 이 식을 외우기보다, 속도 관계보다 가속도 관계가 한 단계 더 까다롭다는 사실만 붙잡아도 좋다.
로봇이 빠르게 움직이고 그 구간에서 \(\bm{J}\)가 자세에 따라 많이 변할수록 \(\dot{\bm{J}}\dot{\bm{q}}\) 같은 항을 무시하기 어려워진다.

3차원 로봇에서는 손끝의 선속도뿐 아니라 각속도도 중요하다.
선속도와 각속도를 함께 묶은 6차원 속도를 트위스트(twist)라고 부르고, 힘과 모멘트를 함께 묶은 6차원 힘을 렌치(wrench)라고 부른다.
이 절의 \(\bm{J}\)는 먼저 병진 위치를 다루는 자코비안 \(\bm{J}_v\)로 읽으면 된다.
실제 6차원 자세 제어에서는 기하학적 자코비안 \(\bm{J}_g\), 트위스트 \(\bm{V}\), 렌치 \(\bm{w}\)를 같은 좌표계와 같은 기준점에서 표현해
\[
  \bm{V}=\bm{J}_g\dot{\bm{q}},
  \qquad
  \bm{\tau}=\bm{J}_g^{\mathsf{T}}\bm{w}
\]
처럼 확장한다.

\subsection{특이점과 조작성: 왜 어떤 자세는 어렵게 느껴지는가}
\label{sec:robotics-singularity-conditioning}

자코비안은 현재 자세에서 가능한 손끝 속도 방향을 보여준다.
그런데 어떤 자세에서는 자코비안의 열들이 서로 비슷한 방향을 가리키거나, 필요한 방향을 충분히 만들지 못한다.
이때 로봇은 특정 방향으로 손끝을 움직이기 어려워진다.
이런 자세를 특이점(singularity)에 가깝다고 말한다.

2링크 팔을 거의 일직선으로 쭉 편 모습을 떠올려 보자.
어깨와 팔꿈치를 조금 돌렸을 때 손끝이 움직일 수 있는 방향들이 서로 비슷해진다.
그러면 어떤 방향의 손끝 속도를 만들기 위해 관절 속도를 매우 크게 요구하거나, 아예 만들기 어려운 방향이 생긴다.
수학적으로는 자코비안의 랭크가 떨어지고, 가장 작은 특이값 \(\sigma_{\min}\)이 0에 가까워지며, 조건수 \(\kappa=\sigma_{\max}/\sigma_{\min}\)가 커지는 현상으로 나타난다.

조작성(manipulability)은 현재 자세에서 손끝을 여러 방향으로 얼마나 쉽게 움직일 수 있는지를 나타내는 지표이다.
값 하나만으로 모든 것을 판단할 수는 없지만, Lab03처럼 교육용 플롯을 볼 때는 조작성이 낮아지고 조건수가 커지는 구간을 특이점 근처의 신호로 읽을 수 있다.
2링크 평면 팔에서는
\begin{equation}
  w
  =
  \sqrt{\det(\bm{J}\bm{J}^{\mathsf{T}})}
  =
  |l_1l_2\sin q_2|
  \label{eq:two-link-manipulability}
\end{equation}
처럼 조작성을 읽을 수 있고, \(q_2=0\) 또는 \(\pi\) 근처에서 팔이 일직선에 가까워지면 \(w\)가 0에 가까워진다.
감쇠 최소제곱(DLS)은 이런 구간에서 역기구학 계산이 너무 큰 관절 속도를 요구하지 않도록 완화하는 방법이다 \cite{buss2005sdls}.
여기서 입력은 원하는 손끝 속도 \(\dot{\bm{x}}_d\)이고, 출력은 그 속도를 만들기 위한 관절 속도 \(\dot{\bm{q}}\)이다.
대표적인 속도 IK 형태로 쓰면
\begin{equation}
  \dot{\bm{q}}
  =
  \bm{J}^{\mathsf{T}}
  \left(\bm{J}\bm{J}^{\mathsf{T}}+\lambda^2\bm{I}\right)^{-1}
  \dot{\bm{x}}_d
  \label{eq:dls-velocity-ik}
\end{equation}
처럼 쓸 수 있다.
이 식은 손끝 힘을 관절 effort로 옮기는 뒤 절의 \(\bm{J}^{\mathsf{T}}\bm{f}\)와 모양이 비슷하지만, 용도는 다르다.
\vmark{effort-dls-force-handoff}%
여기서는 원하는 손끝 속도 \(\dot{\bm{x}}_d\)를 관절 속도 \(\dot{\bm{q}}\)로 바꾸는 속도 IK 식으로 읽어야 한다.
한 번에 외우려 하지 말고 오른쪽에서 왼쪽으로 읽으면 덜 어렵다.

\begin{table}[!htbp]
  \centering
  \caption{DLS 속도 IK 식을 오른쪽에서 왼쪽으로 읽는 법}
  \label{tab:dls-velocity-ik-reading}
  \small
  \renewcommand{\arraystretch}{1.12}
  \begin{tabularx}{\linewidth}{@{}p{0.28\linewidth}X@{}}
    \toprule
    식의 조각 & 초보자용 읽기 \\
    \midrule
    \(\dot{\bm{x}}_d\) & 지금 만들고 싶은 손끝 속도이다 \\
    \(\bm{J}\bm{J}^{\mathsf{T}}\) & 현재 자세에서 손끝 속도를 만들기 쉬운 방향과 어려운 방향을 담은 작업공간 쪽 행렬이다 \\
    \(\lambda^2\bm{I}\) & 너무 얇아진 방향을 수치적으로 조금 두껍게 만들어 역계산이 폭주하지 않게 한다 \\
    \((\bm{J}\bm{J}^{\mathsf{T}}+\lambda^2\bm{I})^{-1}\) & 원하는 손끝 속도 중 위험하게 증폭될 부분을 완화해서 읽는다 \\
    \(\bm{J}^{\mathsf{T}}\) & 완화된 작업공간 속도 요구를 관절 속도 방향으로 되돌린다 \\
    \bottomrule
  \end{tabularx}
\end{table}

여기서 \(\lambda\)는 물리 댐퍼가 아니라 수치 계산이 과하게 커지지 않도록 넣는 완화 파라미터이다.
\(\bm{I}\)는 행렬 크기를 맞추기 위한 단위행렬이고, \(\lambda\)를 키우면 움직임은 더 보수적이지만 관절 속도 폭주는 덜해진다.
특이값 관점에서는 순수 역계산이 작은 특이값 \(\sigma\) 방향에서 대략 \(1/\sigma\)처럼 커지는 반면, DLS는 그 방향의 이득을
\begin{equation}
  g_{\mathrm{DLS}}(\sigma)
  =
  \frac{\sigma}{\sigma^2+\lambda^2}
  \label{eq:dls-singular-value-gain}
\end{equation}
처럼 바꾸어 분모를 \(\lambda^2\)로 받쳐 준다.
예를 들어 \(\sigma=0.05\), \(\lambda=0.1\)이면 순수 역계산의 이득 \(1/\sigma=20\) 대신 DLS 이득은 \(0.05/(0.05^2+0.1^2)=4\)가 된다.
관절 속도 요구는 줄어들지만, 실제 \(\bm{J}\dot{\bm{q}}\)가 \(\dot{\bm{x}}_d\)를 완전히 따라가지 못하는 잔차가 커질 수 있다.
즉 DLS는 정확도를 조금 양보하고 관절 속도 폭주를 줄이는 타협이다.
현재 본문의 2차원 병진 속도 IK에서는 이 설명으로 충분하지만, 6차원 트위스트까지 확장할 때는 선속도와 각속도의 단위, 좌표계, 기준점, 스케일링을 함께 맞추어야 한다.

\paragraph{체크포인트.}
\vmark{nav-jacobian-checkpoint}%
여기까지의 \(\bm{J}\)는 현재 자세에서의 국소 미분이다.
자코비안 열은 해당 관절만 단위 속도로 움직일 때의 손끝 순간속도로 읽고, \(\bm{J}\dot{\bm{q}}\)는 관절 속도를 손끝 속도로 옮기는 속도 번역으로 읽는다.
DLS 식 전체는 원하는 손끝 속도를 관절 속도로 바꾸는 속도 IK 완화식이다.
\vmark{nav-dls-velocity-ik-relief}%
다음 절의 \(\bm{J}^{\mathsf{T}}\bm{f}\)는 모양은 비슷하지만 일률 또는 가상일에서 나온 힘--관절 effort 관계이므로, DLS의 속도 역계산과 섞어 읽지 않는다.
\vmark{nav-dls-jtranspose-distinction}%
또 이 힘 관계를 쓸 때는 \(\bm{f}\), \(\dot{\bm{x}}\), \(\bm{J}\)의 좌표계와 힘의 부호 약속을 함께 확인해야 한다.
\vmark{nav-same-frame-sign}%

\subsection{힘과 토크: 왜 \texorpdfstring{\(\bm{J}^{\mathsf{T}}\)}{J transpose}가 나오는가}
\label{sec:robotics-force-torque-jacobian-transpose}

자코비안은 속도만 연결하지 않는다.
손끝 힘과 관절 effort도 연결한다.
\vmark{effort-force-intro}%
이때는 자코비안 자체가 아니라 전치 \(\bm{J}^{\mathsf{T}}\)가 등장한다 \cite{lynchpark2017modernrobotics}.

출발점은 일률(power)이다.
여기서 \(\bm{f}\)는 선택한 작업공간 방향에서 로봇이 만들려고 하는 힘 명령으로 두자.
\vmark{nav-friendly-task-space-direction}%
손끝에서 힘 \(\bm{f}\)가 손끝 속도 \(\dot{\bm{x}}\)와 같은 방향으로 작용하면 순간 일률은
\begin{equation}
  P_x
  =
  \bm{f}^{\mathsf{T}}\dot{\bm{x}}
\end{equation}
이다.
관절 쪽에서 일반화 관절 effort \(\bm{\tau}\)가 관절 속도 \(\dot{\bm{q}}\)에 대해 하는 일률은
\vmark{effort-generalized-power}%
\begin{equation}
  P_q
  =
  \bm{\tau}^{\mathsf{T}}\dot{\bm{q}}
\end{equation}
이다.
마찰이나 손실을 잠시 빼고 같은 순간의 힘 전달만 보면, 손끝에서 본 일률과 관절에서 본 일률은 같은 물리적 일을 다른 좌표로 쓴 것이다.
따라서
\begin{equation}
  \bm{f}^{\mathsf{T}}\dot{\bm{x}}
  =
  \bm{\tau}^{\mathsf{T}}\dot{\bm{q}}
\end{equation}
라고 둔다.

여기에 속도 관계 \(\dot{\bm{x}}=\bm{J}\dot{\bm{q}}\)를 넣으면
\begin{equation}
  \bm{f}^{\mathsf{T}}\bm{J}\dot{\bm{q}}
  =
  \bm{\tau}^{\mathsf{T}}\dot{\bm{q}}
\end{equation}
이다.
왼쪽은 행렬 전치 규칙으로
\[
  \bm{f}^{\mathsf{T}}\bm{J}\dot{\bm{q}}
  =
  (\bm{J}^{\mathsf{T}}\bm{f})^{\mathsf{T}}\dot{\bm{q}}
\]
처럼 쓸 수 있다.
이 관계가 어떤 관절 속도 \(\dot{\bm{q}}\)에 대해서도 맞으려면
\begin{equation}
  \bm{\tau}
  =
  \bm{J}(\bm{q})^{\mathsf{T}}\bm{f}
  \label{eq:jacobian-transpose-foundation}
\end{equation}
가 된다.
같은 이야기는 아주 작은 가상 변위로도 읽을 수 있다.
\(\delta\bm{x}=\bm{J}\delta\bm{q}\)이면
\[
  \bm{f}^{\mathsf{T}}\delta\bm{x}
  =
  (\bm{J}^{\mathsf{T}}\bm{f})^{\mathsf{T}}\delta\bm{q}
\]
가 되므로, \(\bm{J}^{\mathsf{T}}\bm{f}\)는 작업공간에서 한 일을 관절공간의 일로 똑같이 읽게 해 주는 일반화 관절 effort이다.

이 식의 의미는 손끝 힘을 만들기 위해 각 관절이 얼마만큼의 일반화 effort를 나눠 맡아야 하는지를 계산한다는 것이다.
\vmark{effort-sharing}%
회전 관절에서는 \(\tau_i\)가 토크이고, 직선 관절에서는 해당 축 방향 힘으로 읽는다.
좀 더 단위로 읽어 보면 다음과 같다.
회전 관절의 자코비안 열은 손끝까지의 순간 팔길이, 즉 모멘트 암(moment arm)을 담고 있어서 대략 \(\mathrm{m/rad}\)로 읽힌다.
여기에 손끝 힘 \(\mathrm{N}\)을 곱하면 \(\mathrm{N\,m/rad}\)가 되고, 라디안을 각도 표식으로 남겨 두면 결국 회전축에 필요한 토크 \(\mathrm{N\,m}\)로 읽는다.
직선 관절의 자코비안 열은 축 방향으로 1 m 밀었을 때 손끝이 어느 방향으로 1 m 움직이는지를 나타내므로 \(\mathrm{m/m}\)이다.
여기에 힘 \(\mathrm{N}\)을 곱하면 그 축을 따라 필요한 선형 힘 \(\mathrm{N}\)이 된다.
\begin{equation}
  \begin{aligned}
    \text{회전 관절:}\quad
    &\tau_i
    =
    \bm{j}_i^{\mathsf{T}}\bm{f}
    \sim
    (\mathrm{m/rad})\cdot\mathrm{N}
    \rightarrow
    \mathrm{N\,m},
    \\
    \text{직선 관절:}\quad
    &\tau_i
    =
    \bm{j}_i^{\mathsf{T}}\bm{f}
    \sim
    (\mathrm{m/m})\cdot\mathrm{N}
    \rightarrow
    \mathrm{N}.
  \end{aligned}
  \label{eq:joint-effort-unit-check}
\end{equation}
예를 들어 손끝에 수직 방향 \(10~\mathrm{N}\)의 힘을 만들고, 어떤 회전 관절의 모멘트 암이 \(0.4~\mathrm{m}\)라면 그 관절은 대략 \(4~\mathrm{N\,m}\)의 토크를 맡는다.
반대로 직선 관절 축이 힘 방향과 정확히 같다면 같은 \(10~\mathrm{N}\)은 그대로 \(10~\mathrm{N}\)의 축 방향 힘으로 읽힌다.
가장 단순한 1차원 직선 관절로 검산하면 더 쉽다.
\vmark{prism-one-dimensional-slider}%
슬라이더 축을 \(x\) 방향으로 잡고 스칼라 좌표를 \(x=q\)로 두면 \(\bm{J}=[1]\), \(\dot{x}=\dot{q}\), \(\tau=\bm{J}^{\mathsf{T}}f=f\)가 된다.
\vmark{prism-scalar-j-equals-one}%
\vmark{prism-scalar-tau-equals-force}%
\vmark{prism-scalar-x-equals-q}%
따라서 이때의 \(\tau\)는 토크가 아니라 \(q\)와 짝을 이루는 축 방향 선형 힘이다.
3차원 힘으로 쓰면 \(+x\) 방향 슬라이더의 자코비안 열은 \(\bm{j}=[1\ 0\ 0]^{\mathsf{T}}\)이고,
\vmark{prism-three-dimensional-axis}%
\[
  \tau
  =
  \bm{j}^{\mathsf{T}}\bm{f}
  =
  f_x
\vmark{prism-projection-fx}%
\]
이다.
즉 축에 수직인 힘 성분은 사라지는 것이 아니라, 그 이상적인 슬라이더 자유도에 일을 하지 않기 때문에 actuator의 일반화 effort로 투영되지 않는다.
\vmark{prism-perpendicular-projection}%
단, \(\bm{f}\), \(\dot{\bm{x}}\), \(\bm{J}\)는 같은 좌표계에서 표현되어야 한다.
손끝 좌표계에서 측정한 힘을 월드 좌표계 자코비안에 바로 넣으면 부호와 방향이 틀어질 수 있다.
또 환경이 로봇에 가하는 외력 \(\bm{f}_{\mathrm{ext}}\)을 그대로 쓰는 문맥에서는 부호가 달라진다.
정적 평형을 단순히 쓰면 \(\bm{\tau}_{\mathrm{cmd}}+\bm{J}^{\mathsf{T}}\bm{f}_{\mathrm{ext}}\approx\bm{0}\)처럼 읽을 수 있다.
따라서 \(\bm{J}^{\mathsf{T}}\bm{f}\)를 볼 때는 이 \(\bm{f}\)가 로봇이 만들려는 힘인지, 환경이 로봇에 가하는 힘인지 먼저 확인해야 한다.
또 식 \eqref{eq:jacobian-transpose-foundation}은 이상적인 순간 관계이지, 실제 구동기가 그 힘을 항상 정확히 만들 수 있다는 보증은 아니다.
구동기 effort 제한, 마찰, 제어 대역폭, 특이점, 모델 오차가 모두 영향을 준다.
\vmark{effort-actuator-limit}%

\(\bm{J}\)의 열을 \(\bm{j}_1,\bm{j}_2,\ldots,\bm{j}_n\)이라고 쓰면 이 식은
\[
  \tau_i
  =
  \bm{j}_i^{\mathsf{T}}\bm{f}
\]
라고 읽을 수 있다.
즉 \(i\)번째 관절 effort는 \(i\)번째 관절이 만들 수 있는 손끝 속도 방향과 원하는 손끝 힘이 얼마나 같은 방향을 보는지에 의해 정해진다.
\vmark{prism-joint-effort-reading}%
회전 관절이면 이 effort가 토크이고, 직선 관절이면 축 방향 힘이다.

\subsection{경로와 궤적: 어디로 갈지와 언제 갈지}
\label{sec:robotics-path-trajectory}

로봇이 움직인다는 말도 두 단계로 나누어 읽어야 한다.
첫째는 어디를 지나갈 것인가이고, 둘째는 그 길을 시간에 따라 얼마나 빠르게 지나갈 것인가이다.
Modern Robotics에서는 이 둘을 경로(path)와 궤적(trajectory)으로 구분한다 \cite{lynchpark2017modernrobotics}.
지도 위에 그은 선이 경로라면, 그 선을 몇 초에 어느 지점까지 갈지 적어 둔 시간표가 궤적이다.
같은 직선 \(10~\mathrm{cm}\)를 움직이더라도 5초에 걸쳐 천천히 가는 경우와 1초 안에 가는 경우는 같은 경로를 지난다.
하지만 속도와 가속도는 다르고, 접촉 직전에 빠르게 접근하면 같은 강성에서도 힘 피크와 진동이 더 크게 보일 수 있다.

경로는 시간표가 빠진 기하학적 길이다.
보통 \(s\in[0,1]\)이라는 경로 파라미터를 두고
\begin{equation}
  \bm{q}=\bm{q}(s)
  \quad \text{또는} \quad
  \bm{x}=\bm{x}(s)
\end{equation}
처럼 쓴다.
\(s=0\)이면 시작점, \(s=1\)이면 끝점이다.
이때는 아직 몇 초에 어디에 있을지가 정해지지 않았다.

궤적은 경로에 시간을 붙인 것이다.
시간 스케일링(time scaling) \(s(t)\)를 정하면
\begin{equation}
  \bm{q}(t)=\bm{q}(s(t))
  \label{eq:trajectory-time-scaling-foundation}
\end{equation}
이 된다.
즉 같은 경로라도 \(s(t)\)를 천천히 바꾸면 느린 궤적이 되고, 빠르게 바꾸면 빠른 궤적이 된다.

속도와 가속도는 연쇄법칙으로 연결된다.
\begin{align}
  \dot{\bm{q}}(t)
  &=
  \frac{\dd \bm{q}}{\dd s}\dot{s}(t),
  \\
  \ddot{\bm{q}}(t)
  &=
  \frac{\dd \bm{q}}{\dd s}\ddot{s}(t)
  +
  \frac{\dd^2 \bm{q}}{\dd s^2}\dot{s}(t)^2.
\end{align}
여기서 초보자가 가장 자주 놓치는 지점은 \(\bm{q}(s)\)나 \(\bm{x}(s)\)가 부드러워 보여도, 시간표 \(s(t)\)를 급하게 만들면 실제 속도와 가속도 요구가 커진다는 점이다.
가장 단순한 예로 손끝이 직선으로 \(10~\mathrm{cm}\) 움직인다고 하자.
\begin{equation}
  \bm{x}(s)
  =
  \bm{x}_0+s\Delta\bm{x},
  \qquad
  \|\Delta\bm{x}\|=0.10~\mathrm{m},
  \qquad
  0\le s\le 1.
  \label{eq:straight-path-time-scaling-checkpoint}
\end{equation}
이 식은 어디를 지나가는지만 정한 경로이다.
시간표를 붙이면 \(\bm{x}(t)=\bm{x}(s(t))\)가 되고,
\begin{equation}
  \dot{\bm{x}}
  =
  \Delta\bm{x}\dot{s},
  \qquad
  \ddot{\bm{x}}
  =
  \Delta\bm{x}\ddot{s},
  \qquad
  \dddot{\bm{x}}
  =
  \Delta\bm{x}\dddot{s}
  \label{eq:straight-path-time-derivatives}
\end{equation}
가 된다.
경로가 직선이라 \(\dd^2\bm{x}/\dd s^2=\bm{0}\)인 특별히 쉬운 경우에도, 시간표의 \(\dot{s}\), \(\ddot{s}\), \(\dddot{s}\)가 그대로 실제 속도, 가속도, 저크 요구로 들어온다.

같은 모양의 시간 스케일링을 단지 전체 시간 \(T\) 안에 압축한다고 생각해 보자.
즉 \(s(t)=\hat{s}(t/T)\)라 두면
\[
  \dot{s}\propto \frac{1}{T},
  \qquad
  \ddot{s}\propto \frac{1}{T^2},
  \qquad
  \dddot{s}\propto \frac{1}{T^3}
\]
처럼 커진다.
따라서 같은 \(10~\mathrm{cm}\) 경로라도 \(5~\mathrm{s}\)짜리 움직임을 \(1~\mathrm{s}\)로 압축하면, 같은 형태의 시간표에서는 대표 속도 요구가 약 \(5\)배, 가속도 요구가 약 \(25\)배, 저크 요구가 약 \(125\)배까지 커질 수 있다.
이 숫자는 특정 로봇의 정확한 힘을 예측하기 위한 공식이라기보다, 왜 같은 경로가 전혀 다른 effort와 접촉감을 만들 수 있는지 보여 주는 계산 체크포인트이다.

\begin{table}[!htbp]
  \centering
  \caption{같은 경로를 다른 시간표로 움직일 때 먼저 보는 차이}
  \label{tab:same-path-different-timing}
  \small
  \setlength{\tabcolsep}{4pt}
  \renewcommand{\arraystretch}{1.12}
  \begin{tabularx}{\linewidth}{@{}>{\raggedright\arraybackslash}p{0.20\linewidth}>{\raggedright\arraybackslash}p{0.30\linewidth}>{\raggedright\arraybackslash}X@{}}
    \toprule
    비교 & 같은 점 & 달라지는 점 \\
    \midrule
    느린 \(5~\mathrm{s}\) 이동 & \(\|\Delta\bm{x}\|=0.10~\mathrm{m}\)인 같은 직선 경로 & 대표 속도, 가속도, 저크 요구가 작아 접촉 직전 과도응답과 effort 피크가 완만해지기 쉽다 \\
    빠른 \(1~\mathrm{s}\) 이동 & 시작점과 끝점, 지나가는 공간 길은 같다 & 시간표가 압축되어 속도와 가속도 요구가 커지고, 벽에 접근할 때 감쇠력과 effort 피크가 커질 수 있다 \\
    접촉 실험에서의 질문 & 경로만 보면 둘 다 같은 목표로 간다 & 접촉 직전 속도, 가속도, 저크가 벽 힘, 튕김, 목표 후퇴, actuator effort를 어떻게 바꾸는지 함께 봐야 한다 \\
    \bottomrule
  \end{tabularx}
\end{table}

따라서 궤적을 설계할 때는 위치만 부드러우면 끝나는 것이 아니라, 속도와 가속도도 구동기가 감당할 수 있는지 확인해야 한다.
접촉 작업에서는 더 조심해야 한다.
벽에 닿기 직전 속도와 가속도가 크면, 같은 강성이라도 접촉 순간 힘과 진동이 크게 나타날 수 있다.

\subsection{Trapezoidal, S-curve, 그리고 jerk}
\label{sec:robotics-trajectory-profiles}

가장 흔한 점대점(point-to-point) 궤적 중 하나는 사다리꼴 속도 프로파일(trapezoidal velocity profile)이다.
속도 그래프를 그리면 처음에는 일정한 가속도로 속도가 증가하고, 중간에는 일정한 속도로 움직이며, 마지막에는 일정한 감속도로 멈춘다.
그래서 속도 그래프가 사다리꼴처럼 보인다.
여기서 꼭 붙잡아야 할 직관은 속도 그래프 아래 면적이 이동거리라는 점이다.
예를 들어 \(L=0.10~\mathrm{m}\)를 움직이고, 최고 속도 \(v_{\max}=0.05~\mathrm{m/s}\), 최대 가속도 \(a_{\max}=0.10~\mathrm{m/s^2}\)로 제한한다고 하자.
가속에 걸리는 시간은
\begin{equation}
  t_{\mathrm{acc}}
  =
  \frac{v_{\max}}{a_{\max}}
  =
  \frac{0.05}{0.10}
  =
  0.5~\mathrm{s}
  \label{eq:trapezoid-accel-time}
\end{equation}
이고, 이때 가속 구간에서 이동한 거리는
\begin{equation}
  d_{\mathrm{acc}}
  =
  \frac{1}{2}a_{\max}t_{\mathrm{acc}}^2
  =
  \frac{1}{2}\cdot0.10\cdot0.5^2
  =
  0.0125~\mathrm{m}
  \label{eq:trapezoid-accel-distance}
\end{equation}
이다.
가속과 감속 구간을 합쳐도 \(0.025~\mathrm{m}\)만 쓰므로, 남은 \(0.075~\mathrm{m}\)는 최고 속도로 움직이는 등속 구간이 된다.
따라서
\begin{equation}
  t_{\mathrm{flat}}
  =
  \frac{L-2d_{\mathrm{acc}}}{v_{\max}}
  =
  \frac{0.10-2\cdot0.0125}{0.05}
  =
  1.5~\mathrm{s},
  \qquad
  T
  =
  2t_{\mathrm{acc}}+t_{\mathrm{flat}}
  =
  2.5~\mathrm{s}.
  \label{eq:trapezoid-total-time}
\end{equation}
이 예제는 사다리꼴 속도 프로파일이 단순히 예쁜 그래프 이름이 아니라, 속도 그래프의 면적으로 이동거리를 맞추는 계산이라는 점을 보여 준다.
만약 \(L<2d_{\mathrm{acc}}\)이면 최고 속도에 도달하기 전에 감속해야 하므로, 속도 그래프는 사다리꼴이 아니라 삼각형(triangular profile)이 된다.

이 방법은 이해하기 쉽고 구현도 단순하다.
하지만 가속도는 구간 경계에서 갑자기 바뀐다.
저크는 속도가 얼마나 빠른지가 아니라, 가속도가 얼마나 갑자기 바뀌는지를 보는 값이다.
그래서 같은 최고 속도에 도달하더라도 가속도를 단번에 켜면 저크가 크고, 가속도를 서서히 올리면 저크가 작다.
이상적인 사다리꼴 속도 프로파일에서는 가속도가 불연속으로 점프하므로, 저크는 전환 순간에 임펄스처럼 나타난다.
실제 로봇에서는 필터, 모터 대역폭, 제어 주기 때문에 무한대가 되지는 않지만, 여전히 큰 저크로 완화되어 진동과 effort 피크를 만들 수 있다.
강체 하나를 아주 천천히 움직일 때는 큰 문제가 아닐 수 있지만, 유연한 구조물, 무거운 부하, 빠른 산업용 로봇, 접촉 직전의 로봇에서는 이런 급격한 변화가 진동과 effort 피크를 만들 수 있다.

S-curve 궤적은 이 문제를 줄이기 위해 가속도 자체도 부드럽게 바꾸려는 방식이다.
전형적인 S-curve는 저크의 크기를 제한하여 가속도가 서서히 올라가고, 필요하면 일정 가속 구간을 지나고, 다시 서서히 내려가게 만든다.
간단히 첫 구간만 보면, 최대 저크를 \(j_{\max}=0.4~\mathrm{m/s^3}\)로 제한하고 \(a_{\max}=0.10~\mathrm{m/s^2}\)까지 가속도를 올리고 싶을 때 걸리는 시간은
\begin{equation}
  t_j
  =
  \frac{a_{\max}}{j_{\max}}
  =
  \frac{0.10}{0.4}
  =
  0.25~\mathrm{s}
  \label{eq:s-curve-jerk-ramp-time}
\end{equation}
이다.
이 \(0.25~\mathrm{s}\) 동안 가속도는 갑자기 \(0.10~\mathrm{m/s^2}\)로 튀는 대신 \(0\)에서 \(0.10~\mathrm{m/s^2}\)까지 선형으로 올라간다.
그 구간에서 속도가 늘어나는 양은 가속도-시간 그래프의 삼각형 면적이므로
\begin{equation}
  \Delta v_j
  =
  \frac{1}{2}j_{\max}t_j^2
  =
  \frac{1}{2}\cdot0.4\cdot0.25^2
  =
  0.0125~\mathrm{m/s}.
  \label{eq:s-curve-jerk-ramp-dv}
\end{equation}
이고, 같은 구간에서 위치가 늘어나는 양은
\begin{equation}
  \Delta x_j
  =
  \frac{1}{6}j_{\max}t_j^3
  =
  \frac{1}{6}\cdot0.4\cdot0.25^3
  \approx
  0.00104~\mathrm{m}
  \label{eq:s-curve-jerk-ramp-dx}
\end{equation}
이다.
이 작은 계산은 전체 S-curve를 다 설계했다는 뜻이 아니라, S-curve가 왜 더 부드럽게 느껴지는지를 보여 주는 첫 단추이다.
Modern Robotics의 궤적 생성 장도 다항식 시간 스케일링, 사다리꼴, S-curve 같은 시간 스케일링을 경로를 궤적으로 바꾸는 방법으로 소개한다 \cite{lynchpark2017modernrobotics}.
최근에는 여러 자유도에 대해 속도, 가속도, 저크 제한을 만족시키며 온라인으로 궤적을 생성하는 방법도 활발히 쓰이며, Ruckig는 그런 jerk-limited online trajectory generation의 대표적인 예로 소개된다 \cite{berscheid2021ruckig}.

\paragraph{프로파일 선택 체크포인트.}
\vmark{traj-profile-selection-checkpoint}%
처음 고를 때는 사다리꼴을 기본값으로 두고, 속도와 가속도 제한만으로 충분한지 먼저 확인하면 된다.
\vmark{traj-trapezoid-default}%
전환 순간의 진동, 소음, 접촉 충격, 무거운 부하가 문제가 되면 그때 저크를 제한하거나 더 완만하게 만드는 S-curve 계열 프로파일을 고려한다.
\vmark{traj-transition-shock}%
다만 S-curve라는 이름이 항상 모든 속도, 가속도, 저크 제한을 엄밀히 만족하는 온라인 생성기라는 뜻은 아니다.
\vmark{traj-scurve-not-always-generator}%
교육용 실험에서는 S-curve-shaped smooth polynomial target을 비교 신호로 쓰기도 하므로, 로그의 jerk나 effort는 같은 시뮬레이터와 컨트롤러 안에서 읽는 진단 신호이지 실제 접촉력이나 하드웨어 성능 검증값은 아니다.
\vmark{traj-diagnostic-not-hardware}%
\vmark{traj-smooth-polynomial-target}%

다음 그림에서는 먼저 같은 경로가 유지되는지 보고, 그다음 속도와 가속도 모양이 어떻게 달라지는지 비교하면 된다.
읽는 순서는 단순하다.
먼저 속도 그래프의 면적이 이동거리라는 점을 확인하고, 다음으로 가속도 그래프가 계단처럼 끊기는지, 마지막으로 저크가 큰 전환점을 갖는지 본다.

\begin{figure}[!htbp]
  \centering
  \resizebox{\linewidth}{!}{%
  \begin{tikzpicture}[
      axis/.style={-{Latex[length=1.7mm]}, draw=black!65, line width=0.45pt},
      curve/.style={line width=0.95pt},
      guide/.style={densely dashed, draw=black!35, line width=0.4pt},
      panel/.style={draw=black!25, rounded corners=2pt, fill=black!2},
      title/.style={font=\scriptsize\bfseries, align=center},
      lab/.style={font=\scriptsize, align=center},
      note/.style={draw=black!35, rounded corners=2pt, fill=white, align=center, font=\scriptsize, inner sep=3.5pt}
    ]

    \begin{scope}[shift={(0,0)}]
      \draw[panel] (-0.25,-0.25) rectangle (3.45,2.85);
      \node[title] at (1.6,2.62) {경로: 어디를 지날까};
      \draw[axis] (0,0) -- (3.05,0) node[right, lab] {$x$};
      \draw[axis] (0,0) -- (0,2.15) node[above, lab] {$y$};
      \draw[curve, blue!70!black] plot[smooth] coordinates {
        (0.35,0.35) (0.85,1.1) (1.55,1.55) (2.25,1.25) (2.8,1.85)
      };
      \fill[blue!70!black] (0.35,0.35) circle (0.055);
      \fill[blue!70!black] (2.8,1.85) circle (0.055);
      \node[lab, anchor=east] at (0.28,0.35) {$s=0$};
      \node[lab, anchor=west] at (2.88,1.85) {$s=1$};
      \node[note] at (1.65,0.12) {지도 위 선\\시간표 없음};
    \end{scope}

    \begin{scope}[shift={(4.05,0)}]
      \draw[panel] (-0.25,-0.25) rectangle (3.45,2.85);
      \node[title] at (1.6,2.62) {궤적: 언제 어디까지 갈까};
      \draw[axis] (0,0) -- (3.05,0) node[right, lab] {$t$};
      \draw[axis] (0,0) -- (0,2.15) node[above, lab] {$s(t)$};
      \draw[guide] (0,1.9) -- (2.95,1.9);
      \node[lab, anchor=east] at (-0.03,1.9) {$1$};
      \draw[curve, green!45!black] plot[smooth] coordinates {
        (0.0,0.0) (0.4,0.35) (0.8,0.92) (1.2,1.4) (1.65,1.72) (2.1,1.86) (2.7,1.9)
      };
      \draw[curve, orange!80!black, dashed] plot[smooth] coordinates {
        (0.0,0.0) (0.65,0.18) (1.15,0.55) (1.7,1.08) (2.25,1.55) (2.75,1.82) (2.95,1.9)
      };
      \node[lab, text=green!45!black] at (1.1,1.85) {빠른 시간표};
      \node[lab, text=orange!80!black] at (2.15,0.7) {느린 시간표};
      \node[note] at (1.55,0.12) {같은 경로라도\\속도와 접촉감이 달라짐};
    \end{scope}

    \begin{scope}[shift={(8.1,1.45)}]
      \draw[panel] (-0.25,-1.7) rectangle (4.45,1.4);
      \node[title] at (2.1,1.18) {사다리꼴 속도};
      \draw[axis] (0,0) -- (4.05,0) node[right, lab] {$t$};
      \draw[axis] (0,0) -- (0,0.95) node[above, lab] {$v$};
      \draw[curve, red!70!black] (0,0) -- (1.0,0.75) -- (2.65,0.75) -- (3.65,0);
      \draw[axis] (0,-1.15) -- (4.05,-1.15) node[right, lab] {$t$};
      \draw[axis] (0,-1.15) -- (0,-0.25) node[above, lab] {$a$};
      \draw[curve, red!70!black] (0,-0.55) -- (1.0,-0.55) -- (1.0,-1.15) -- (2.65,-1.15) -- (2.65,-1.52) -- (3.65,-1.52);
      \draw[guide] (1.0,-1.6) -- (1.0,0.85);
      \draw[guide] (2.65,-1.6) -- (2.65,0.85);
      \node[lab] at (2.05,-1.42) {가속도 점프\\저크 피크};
    \end{scope}

    \begin{scope}[shift={(13.15,1.45)}]
      \draw[panel] (-0.25,-1.7) rectangle (4.45,1.4);
      \node[title] at (2.1,1.18) {S-curve};
      \draw[axis] (0,0) -- (4.05,0) node[right, lab] {$t$};
      \draw[axis] (0,0) -- (0,0.95) node[above, lab] {$v$};
      \draw[curve, purple!70!black] plot[smooth] coordinates {
        (0,0) (0.35,0.08) (0.75,0.38) (1.25,0.68) (1.9,0.75) (2.55,0.68) (3.05,0.38) (3.45,0.08) (3.75,0)
      };
      \draw[axis] (0,-1.15) -- (4.05,-1.15) node[right, lab] {$t$};
      \draw[axis] (0,-1.15) -- (0,-0.25) node[above, lab] {$a$};
      \draw[curve, purple!70!black] plot[smooth] coordinates {
        (0,-1.15) (0.35,-0.75) (0.8,-0.55) (1.2,-0.75) (1.55,-1.15)
        (2.15,-1.15) (2.55,-1.55) (3.0,-1.75) (3.45,-1.55) (3.75,-1.15)
      };
      \node[lab, text=purple!70!black] at (2.05,-1.42) {가속도도\\완만히 변화};
    \end{scope}
  \end{tikzpicture}%
  }
  \caption{경로는 공간에서 지나는 길이고, 궤적은 그 길에 시간표를 붙인 것이다. 같은 경로라도 시간 스케일링 \(s(t)\)이 다르면 속도, 가속도, 접촉 순간의 힘 피크가 달라진다. 사다리꼴 속도 프로파일은 단순하지만 이상적인 형태에서는 가속도가 구간 경계에서 점프하므로 저크가 크게 나타난다. S-curve는 저크 크기를 제한해 가속도 변화를 더 부드럽게 만드는 시간 스케일링이다.}
  \label{fig:path-trajectory-profiles}
\end{figure}


\begin{table}[!htbp]
  \centering
  \caption{경로와 궤적을 처음 구분할 때 보는 기준}
  \label{tab:path-trajectory-profile}
  \small
  \renewcommand{\arraystretch}{1.12}
  \begin{tabularx}{\linewidth}{@{}p{0.18\linewidth}p{0.35\linewidth}X@{}}
    \toprule
    이름 & 뜻 & 플롯에서 보는 것 \\
    \midrule
    경로(path) & 시간표가 없는 기하학적 길 & 손끝이나 관절이 어느 점들을 지나가는가 \\
    궤적(trajectory) & 경로에 시간 스케일링을 붙인 것 & 위치, 속도, 가속도가 시간에 따라 어떻게 변하는가 \\
    사다리꼴 속도 & 가속--등속--감속 구간을 갖는 단순 프로파일 & 속도는 사다리꼴, 가속도는 계단처럼 바뀜 \\
    S-curve & 저크를 제한해 가속도를 부드럽게 바꾸는 프로파일 & 가속도가 갑자기 튀지 않고 완만히 변함 \\
    \bottomrule
  \end{tabularx}
\end{table}

여기서 중요한 것은 사다리꼴이 나쁘고 S-curve가 항상 좋다는 식의 단순한 결론이 아니다.
목표는 로봇과 작업에 맞는 부드러움, 속도, 제한 조건의 균형을 잡는 것이다.
실험실에서는 먼저 단순한 궤적으로 시작해도 된다.
다만 속도, 가속도, 저크가 접촉 힘과 진동에 영향을 줄 수 있다는 감각을 갖고 있어야 한다.

\subsection{다관절 제어와 접촉으로 가는 길}
\label{sec:robotics-multijoint-contact-bridge}

이제 1차원 질량--스프링--댐퍼 직관이 왜 바로 로봇 끝단으로 옮겨지지 않는지 보인다.
1차원에서는 위치 \(x\), 속도 \(\dot{x}\), 힘 \(f\)가 같은 축 위에 있었다.
하지만 다관절 로봇에서는 관절 변수 \(\bm{q}\), 손끝 위치 \(\bm{x}\), 관절 속도 \(\dot{\bm{q}}\), 손끝 속도 \(\dot{\bm{x}}\), 관절 effort \(\bm{\tau}\) (회전 관절에서는 토크), 손끝 힘 \(\bm{f}\)가 서로 연결되어 있다.
\vmark{effort-closing-joint-tau}%
그 연결의 중심에 정기구학과 자코비안이 있다.

자유공간에서는 주로 목표 궤적을 얼마나 잘 따라가는지가 중요하다.
어느 경로를 택할지, 관절공간에서 보간할지 작업공간에서 보간할지, 속도와 가속도 제한을 어떻게 지킬지가 핵심 질문이 된다.
하지만 접촉이 시작되면 질문이 바뀐다.
로봇이 목표 위치를 끝까지 밀고 들어가야 하는가, 아니면 어느 정도 힘이 생기면 부드럽게 물러나야 하는가?
벽 방향은 부드럽게 하고, 벽을 따라 움직이는 방향은 단단하게 유지할 수 있는가?
이 질문은 단순한 위치 추종이 아니라 힘과 운동의 관계를 설계하는 문제이다.

그래서 다음 장의 임피던스 제어가 필요하다.
관절공간과 작업공간의 언어, 자코비안의 속도 변환, \(\bm{J}^{\mathsf{T}}\)의 힘--관절 effort 관계, 궤적의 속도/가속도/저크 감각을 알고 나면, 임피던스 제어는 갑자기 튀어나온 고급 제어가 아니라 자연스러운 다음 질문이 된다.
\vmark{effort-closing-relation}%
흐름만 아주 짧게 쓰면 다음과 같다.
\[
  \text{손끝 오차}
  \rightarrow
  \text{가상 스프링-댐퍼 힘}
  \rightarrow
  \bm{J}^{\mathsf{T}}\text{를 통한 관절 effort}
\vmark{effort-closing-jtf-flow}%
  \rightarrow
  \text{접촉 중 힘과 운동}
\]
여기에 접근 속도, 가속도, 저크가 더해지면 접촉 순간의 튐과 구동기 부담 또는 관절 effort 피크가 달라진다.
\vmark{effort-closing-peak}%
손끝에 어떤 가상 스프링과 댐퍼를 걸고, 그 힘-운동 관계를 실제 관절 명령과 플롯에서 어떻게 읽을 것인가?
이것이 다음 장의 주제이다.
